\chapter{Polynomial Rings and Factorization}

Polynomials turn arithmetic in a ring into a tool for studying roots, maps,
and later linear operators.  Over a field, division of polynomials recovers
the Euclidean algorithm; this leads from fields to PIDs and UFDs.

\section{Polynomial Rings and Evaluation}
For a commutative ring \(R\), \(R[x]\) consists of finite sums
\(a_0+a_1x+\cdots+a_nx^n\), with coefficientwise addition and convolution
multiplication.  If \(R\) is a domain and \(f,g\ne0\), then
\(\deg(fg)=\deg f+\deg g\), because the product of leading coefficients is
nonzero.
\begin{proposition}[Degree Rule]
If \(R\) is a domain and \(f,g\in R[x]\) are nonzero, then
\(\deg(fg)=\deg f+\deg g\).  The units of \(R[x]\) are exactly the constant
polynomials whose value is a unit of \(R\).
\end{proposition}
\begin{proof}
Writing \(f=a_mx^m+\cdots\), \(g=b_nx^n+\cdots\), the coefficient of
\(x^{m+n}\) in \(fg\) is \(a_mb_n\), which is nonzero precisely because a
domain has no zero divisors.  Thus a product can have degree zero only when
both factors have degree zero, proving the assertion about units.
\end{proof}
\begin{proposition}[Evaluation] For \(a\in R\), evaluation
\(\operatorname{ev}_a:R[x]\to R\), \(f\mapsto f(a)\), is a ring
homomorphism.  If \(R\) is a field, its kernel is \((x-a)\). \end{proposition}
\begin{proof} The first assertion follows by distributing finite sums.  The
division algorithm below writes \(f=(x-a)q+r\), where \(r\) is constant;
then \(r=f(a)\), proving the kernel statement. \end{proof}

\section{Division Algorithms and Roots}
\begin{theorem}[Division Algorithm] If \(F\) is a field and \(f,g\in F[x]\),
\(g\ne0\), then there are unique \(q,r\in F[x]\) with
\[f=qg+r,\qquad r=0\ \text{or}\ \deg r<\deg g.\]
\end{theorem}
\begin{proof} Write \(g=b_nx^n+\cdots\), with \(b_n\ne0\).  If the current
dividend is \(h=c_mx^m+\cdots\) with \(m\ge n\), subtract
\((c_mb_n^{-1})x^{m-n}g\).  This is possible because \(F\) is a field, and
it cancels the \(x^m\)-term, leaving a polynomial of smaller degree.  Iterating
must stop, since nonnegative degrees cannot decrease forever; collecting the
subtracted terms gives \(f=qg+r\).  If two such expressions exist, then
\((q-q')g=r'-r\).  Unless both sides are zero, the left side has degree at
least \(\deg g\), while the right side has smaller degree, a contradiction.
\end{proof}
\begin{example}[Long Division]
Dividing \(x^3-2x+4\) by \(x-1\), first subtract \(x^2(x-1)\), then
\(x(x-1)\), then \(-1(x-1)\).  Thus
\[x^3-2x+4=(x^2+x-1)(x-1)+3.\]
\end{example}
\begin{corollary}[Factor Theorem] For \(f\in F[x]\), \(f(a)=0\) iff
\(x-a\) divides \(f\). \end{corollary}
\begin{proof} Apply the division algorithm to \(f\) by \(x-a\) and evaluate.
\end{proof}
\begin{corollary} A nonzero polynomial of degree \(n\) over a field has at
most \(n\) roots. \end{corollary}
\begin{proof} Each root supplies a distinct linear factor; induction on
degree proves the assertion. \end{proof}

\section{Euclidean Domains and Principal Ideal Domains}
\begin{definition} A domain \(R\) is \emph{Euclidean} if it has a function
\(\delta:R\setminus\{0\}\to\mathbb N\) such that division with remainder
reduces \(\delta\).  A \emph{principal ideal domain} (PID) is a domain in
which every ideal is principal. \end{definition}
\begin{definition}
For \(a,b\) in a domain, a \emph{greatest common divisor} is an element \(d\)
such that \(d\mid a\), \(d\mid b\), and every common divisor of \(a,b\)
divides \(d\).  A gcd, when it exists, is unique only up to associates.  In
\(\mathbb Z\) one chooses the positive representative, but a general domain
need not have a preferred associate.
\end{definition}
\begin{theorem} Every Euclidean domain is a PID.  In particular \(F[x]\) is a
PID for every field \(F\). \end{theorem}
\begin{proof} Let \(I\ne0\) be an ideal and choose \(d\in I\) with minimal
\(\delta(d)\).  Divide any \(a\in I\) by \(d\): \(a=qd+r\).  Then
\(r\in I\), so minimality gives \(r=0\), hence \(I=(d)\).  Polynomial degree
is a Euclidean function by the division algorithm. \end{proof}
\begin{proposition}[B\'ezout and gcd]
In a PID, \((a,b)=(d)\) if and only if \(d\) is a gcd of \(a,b\), up to
associates.  Consequently every pair has a gcd and
\[d=ra+sb\]
for suitable \(r,s\).
\end{proposition}
\begin{proof}
If \((a,b)=(d)\), then \(a,b\in(d)\), so \(d\mid a,b\).  A common divisor
\(c\) divides every linear combination of \(a,b\), and hence divides
\(d\in(a,b)\).  Thus \(d\) is a gcd.  Conversely, write
\((a,b)=(g)\), which is possible because \(R\) is a PID.  The forward
direction shows that \(g\) is a gcd.  If \(d\) is another gcd, then
\(d\mid g\), since \(d\) is a common divisor, and \(g\mid d\), since \(g\)
is a common divisor.  Thus \(d\) and \(g\) are associates, so
\((d)=(g)=(a,b)\).  Consequently \(d\in(a,b)\), giving the displayed
B\'ezout representation.  This connects gcds, generated ideals, and
B\'ezout identities.
\end{proof}

\section{Unique Factorization Domains and Gauss's Lemma}
The factorization of an integer into primes is a model for a property enjoyed
by many domains.  The central question here is whether the property survives
the passage from a UFD \(R\) to \(R[x]\).  The answer is yes, and Gauss's
lemma explains why coefficients do not introduce hidden factorizations.

\begin{definition}
Let \(R\) be a domain.  Nonzero \(a,b\in R\) are \emph{associates} if
\(a=ub\) for a unit \(u\).  A nonzero nonunit \(p\) is \emph{irreducible} if
\(p=ab\) implies that \(a\) or \(b\) is a unit.  It is \emph{prime} if
\(p\mid ab\) implies \(p\mid a\) or \(p\mid b\).  A \emph{UFD} is a domain
where every nonzero nonunit has a factorization into irreducibles, unique up
to ordering and replacement by associates.
\end{definition}

Every prime is irreducible: if \(p=ab\), then \(p\mid a\) or \(p\mid b\);
cancellation makes the other factor a unit.  The converse need not hold in
an arbitrary domain, but it does hold in a PID.

\begin{proposition}[Irreducibles are prime in a PID]
Every irreducible element of a principal ideal domain is prime.
\end{proposition}
\begin{proof}
Let \(p\) be irreducible, suppose \(p\mid ab\), and assume \(p\nmid a\).
The gcd \(d\) of \(p,a\) cannot be associated to \(p\), so it is a unit.
Thus \((p,a)=R\), and B\'ezout gives \(rp+sa=1\).  After multiplication by
\(b\), both terms on the left of \(rpb+sab=b\) are divisible by \(p\).
Hence \(p\mid b\).
\end{proof}

\phantomsection\label{thm:irreducible-maximal-field}
\begin{applicationtheorem}[Irreducible polynomials and field quotients]
Let \(F\) be a field and let \(f\in F[x]\) be nonconstant.  The following are
equivalent:
\[
f\text{ is irreducible},
\qquad
(f)\text{ is a maximal ideal of }F[x],
\qquad
F[x]/(f)\text{ is a field}.
\]
\end{applicationtheorem}
\begin{proof}
Since \(F[x]\) is a PID, every ideal containing \((f)\) has the form
\((g)\), where \(g\mid f\).  If \(f\) is irreducible, then \(g\) is either a
unit or an associate of \(f\).  Thus there is no ideal strictly between
\((f)\) and \(F[x]\), so \((f)\) is maximal.  Conversely, if
\(f=gh\) with both \(g\) and \(h\) nonunits, then
\[
(f)\subsetneq(g)\subsetneq F[x],
\]
contradicting maximality.  The equivalence between maximal ideals and field
quotients is the quotient characterization of maximal ideals from Chapter~4.
\end{proof}

Irreducibility therefore does more than forbid a factorization: it is exactly
what makes the quotient obtained by adjoining a formal root into a field.

\begin{theorem}[Every PID is a UFD]
Every principal ideal domain is a unique factorization domain.
\end{theorem}
\begin{proofidea}
Existence comes from the ascending-chain condition on ideals: an element that
could not be factored would force a strictly increasing chain of principal
ideals.  Uniqueness follows by cancelling prime factors.
\end{proofidea}
\begin{proof}
An ascending chain of ideals in a PID stabilizes, since its union is an ideal
\((d)\), and \(d\) belongs to one member of the chain.  Suppose some nonzero
nonunit does not factor into irreducibles, and choose \(a\) with \((a)\)
maximal among ideals generated by such elements.  Then \(a=bc\), with both
\(b,c\) nonunits.  Therefore \((a)\subsetneq(b)\) and
\((a)\subsetneq(c)\), so maximality gives factorizations of \(b,c\), a
contradiction.  Thus factorizations exist.

If \(p_1\cdots p_m=q_1\cdots q_n\) are products of irreducibles, \(p_1\) is
prime by the proposition, hence divides some \(q_j\).  Since \(q_j\) is
irreducible, they are associates.  Reorder, absorb the unit, and cancel.
Induction gives uniqueness.
\end{proof}

Let \(R\) now be a UFD.  A gcd of coefficients is only determined up to a
unit, precisely the ambiguity that factorization theory permits.

Gcds exist in every UFD, although they need not have B\'ezout
representations.  After choosing representatives of irreducible associate
classes, write
\[a=u\prod p_i^{\alpha_i},\qquad b=v\prod p_i^{\beta_i}.\]
Then, up to associates,
\[\gcd(a,b)\sim\prod p_i^{\min(\alpha_i,\beta_i)}.\]
The exponents show at once that this element divides both \(a\) and \(b\),
and that every common divisor divides it.  In a PID, gcds are B\'ezout
combinations; in a general UFD this is false.  For example,
\(\gcd(x,y)=1\) in \(F[x,y]\), but \((x,y)\ne F[x,y]\): an identity
\(fx+gy=1\) would become \(0=1\) after evaluation at \(x=y=0\).
Thus factorization-coprime and comaximal agree in a PID but not in every UFD.

\begin{definition}[Content]
For \(0\ne f=a_0+\cdots+a_nx^n\in R[x]\), its \emph{content} is
\[\operatorname{cont}(f)=\gcd(a_0,\ldots,a_n),\]
where the gcd is understood up to multiplication by a unit.  The polynomial
is \emph{primitive} exactly when \(\operatorname{cont}(f)\) is a unit,
equivalently when no nonunit divides all of its coefficients.
\end{definition}

The content is therefore simply the gcd of the coefficients, whose existence
is supplied by the UFD property just discussed.

Every nonzero \(f\) has the form
\[f=\operatorname{cont}(f)f_0\]
up to multiplication by a unit, with \(f_0\) primitive.
For example,
\[6x^2-9x+3=3(2x^2-3x+1),\]
and the parenthesized factor is primitive in \(\mathbb Z[x]\).

\begin{lemma}[Primitive Product Lemma]
The product of two primitive polynomials in \(R[x]\) is primitive.
\end{lemma}
\begin{proof}
If a nonunit divided every coefficient of \(fg\), some irreducible factor
\(p\) would do so.  In a UFD, irreducibles are prime, so \((p)\) is a prime
ideal: \(ab\in(p)\) means \(p\mid ab\).  Hence \(R/(p)\) is a domain by
Chapter~4.  The reductions of \(f\) and \(g\) modulo \(p\) are nonzero,
because \(f,g\) are primitive, but their product is zero.  This contradicts
the degree rule in the polynomial ring over a domain.
\end{proof}

\begin{corollary}[Content formula]
For nonzero \(f,g\in R[x]\),
\[\operatorname{cont}(fg)\sim\operatorname{cont}(f)\operatorname{cont}(g),\]
where \(\sim\) denotes association.  Equality is naturally only up to
associates in a general UFD.  For \(R=\mathbb Z\), normalizing content to
be positive gives the literal equality
\(\operatorname{cont}(fg)=\operatorname{cont}(f)\operatorname{cont}(g)\).
\end{corollary}
\begin{proof}
Write \(f=\operatorname{cont}(f)f_0\),
\(g=\operatorname{cont}(g)g_0\), where \(f_0,g_0\) are primitive.
Then \(fg=\operatorname{cont}(f)\operatorname{cont}(g)f_0g_0\), and the
lemma says the final factor has unit
content.
\end{proof}

Let \(K=\operatorname{Frac}(R)\).  The next statement is the precise form
of clearing denominators.

\begin{lemma}[Gauss's factorization lemma]
If \(f\in R[x]\) is primitive and \(f=gh\) in \(K[x]\), then there are
primitive \(g_0,h_0\in R[x]\) and a unit \(\varepsilon\in R\) such that
\[f=(\varepsilon g_0)h_0.\]
Thus a positive-degree factorization over \(K\) yields one over \(R\).
\end{lemma}
\begin{proof}
Choose nonzero \(a,b\in R\) such that \(G=ag\) and \(H=bh\) lie in \(R[x]\).
Write \(G=\operatorname{cont}(G)g_0\),
\(H=\operatorname{cont}(H)h_0\) with primitive parts.  Since
\(GH=abf\), the content formula gives
\[\operatorname{cont}(G)\operatorname{cont}(H)
\sim\operatorname{cont}(GH)=\operatorname{cont}(abf)\sim ab.\]
Thus \(\operatorname{cont}(G)\operatorname{cont}(H)=\varepsilon ab\) for a
unit \(\varepsilon\), and
\[abf=GH=\varepsilon abg_0h_0.\]
Cancellation in the domain \(R[x]\) gives the result.  Scaling does not
change degree.
\end{proof}

\begin{corollary}[Irreducibility transfer]
For primitive \(f\in R[x]\),
\[f\text{ is irreducible in }R[x]\quad\Longleftrightarrow\quad
f\text{ is irreducible in }K[x].\]
\end{corollary}
\begin{proof}
A factorization in \(R[x]\) is one in \(K[x]\).  Conversely, a nontrivial
factorization in \(K[x]\) has two positive-degree factors; the lemma turns it
into one in \(R[x]\).
\end{proof}

\begin{lemma}
If primitive \(g,h\in R[x]\) are associates in \(K[x]\), then they are
associates in \(R[x]\).
\end{lemma}
\begin{proof}
Write \(h=(a/b)g\) with \(a,b\in R\setminus\{0\}\).  Then \(bh=ag\).
Taking contents gives \(b\sim a\), so \(a/b\) is a unit of \(R\).
\end{proof}

\begin{theorem}[Polynomial UFD Theorem]
If \(R\) is a UFD, then \(R[x]\) is a UFD.  Hence
\(R[x_1,\ldots,x_n]\) is a UFD for every \(n\ge1\).
\end{theorem}
\begin{proof}
The field \(K\) has \(K[x]\) a PID, hence a UFD.  For \(0\ne f\in R[x]\),
factor its content in \(R\), then factor its primitive part in \(K[x]\).
The factorization lemma clears the denominators of the latter factors; the
irreducibility-transfer corollary makes the resulting primitive factors
irreducible in \(R[x]\).  Irreducible constants of \(R\) remain irreducible
in \(R[x]\), by the degree rule.  This proves existence.

For uniqueness, compare two factorizations in the UFD \(K[x]\).  Its
uniqueness pairs the primitive factors up to association, and the preceding
lemma upgrades these associations to \(R[x]\).  Unique factorization of the
remaining contents occurs in \(R\).  Repeated application of the one-variable
result gives the multivariable assertion.
\end{proof}

\begin{corollary}
\(\mathbb Z[x]\) and \(\mathbb Q[x]\) are UFDs.  A primitive polynomial in
\(\mathbb Z[x]\) is irreducible over \(\mathbb Z\) exactly when it is
irreducible over \(\mathbb Q\).
\end{corollary}
\begin{proof}
Apply the theorem to the UFD \(\mathbb Z\), whose fraction field is
\(\mathbb Q\), and use irreducibility transfer.
\end{proof}

\section{Irreducibility Criteria and Factorization over Fields}
\begin{proposition}
A polynomial of degree two or three over a field is reducible if and only if
it has a root in that field.
\end{proposition}
\begin{proof} A root yields a linear factor by the Factor Theorem.  Conversely,
a proper factorization of degree two or three has a factor of degree one,
which has a root. \end{proof}
\begin{theorem}[Eisenstein's Criterion]
Let \(f=a_nx^n+\cdots+a_0\in\mathbb Z[x]\).  If a prime \(p\) divides every
\(a_i\) for \(i<n\), does not divide \(a_n\), and \(p^2\nmid a_0\), then
\(f\) is irreducible over \(\mathbb Q\).
\end{theorem}
\begin{proof} By Gauss's Lemma it is enough to rule out a factorization in
\(\mathbb Z[x]\); the factors may be chosen primitive.  Hence neither has
all coefficients divisible by \(p\), so their reductions modulo \(p\) are
nonzero.  Reducing a hypothetical \(f=gh\) modulo \(p\) gives
\(\overline g\,\overline h=\overline {a_n}x^n\).  Since the reduced factors
are nonzero and \(\mathbb F_p[x]\) is a domain, each is a nonzero scalar
times a power of \(x\).  Thus every nonleading coefficient of \(g,h\) is
divisible by \(p\).  Their constant terms are therefore both divisible by
\(p\), contradicting \(p^2\nmid a_0\).
\end{proof}
\begin{example} Eisenstein at \(p=2\) proves \(x^n-2\) irreducible over
\(\mathbb Q\) for every \(n\ge1\). \end{example}

\section*{Exercises}
\begin{exercise}[Basic] Divide \(x^3-1\) by \(x-1\) and factor the result over \(\mathbb Q\). \end{exercise}
\begin{exercise}[Core] Use the Euclidean algorithm in \(\mathbb Q[x]\) to
find a gcd of \(x^3-1\) and \(x^2-1\), and express it as a B\'ezout combination.
\end{exercise}
\begin{exercise}[Core] Determine whether \(x^3-2\) and \(x^4+1\) are irreducible over \(\mathbb Q\). \end{exercise}
\begin{exercise}[Core] Find the content and primitive part of each of
\(12x^3-18x^2+6x\) and \(15x^2+10x+5\).  Verify the content formula for
their product. \end{exercise}
\begin{exercise}[Core] Prove directly from the definition that a primitive
polynomial in \(\mathbb Z[x]\) cannot have all of its coefficients divisible
by the same prime.  Use reduction modulo \(2\) to show that
\(x^3+x+1\) is irreducible over \(\mathbb Q\).
\emph{Hint:} use Gauss's Lemma; a nontrivial factorization of this primitive
polynomial would reduce to one in \(\mathbb F_2[x]\), where the cubic has no
root. \end{exercise}
\begin{exercise}[Further] Use B\'ezout to show that coprime polynomials in \(F[x]\) have no common root in any extension field. \end{exercise}
\begin{exercise}[Further, Challenging] Let \(R\) be a UFD and let
\(f\in R[x]\) be primitive.  Show that if \(f\) divides \(g\) in
\(\operatorname{Frac}(R)[x]\) and \(g\in R[x]\), then there is a
primitive polynomial \(h\in R[x]\) and \(a\in R\) such that
\(g=ahf\) up to multiplication by a unit. \end{exercise}
